% titel: Blatt 8, A1: Python-Code
% Teil II: der durchgerechnete Weg.

\subsection*{Gegeben und gesucht}
Den Python-Summenausdruck als mathematische Formel schreiben (a) und die
Newton-Schleife erklären (b).

\subsection*{a) Summenformel}
\begin{verbatim}
sum([1/(k**p) for k in range(1, n+1)])
\end{verbatim}
läuft über $k=1,2,\dots,n$ und summiert $1/k^p$, also (reelle) Summe
\begin{equation*}
  S_n(p)=\sum_{k=1}^{n}\frac{1}{k^p}.
\end{equation*}
Für $p=1$ wächst $S_n(1)\to\infty$ (harmonische Reihe), für $p>1$
konvergiert $S_n(p)\to\zeta(p)$ (Zeta-Funktion).

\subsection*{b) Newton-Schleife}
\begin{verbatim}
x = 3
for n in range(11):
    print(x)
    x = x - (x**3 - 3) / (3*x**2)
\end{verbatim}
Der Code sucht eine Nullstelle von $f(x)=x^3-3$ (also die Kubikwurzel
$\sqrt[3]{3}$). Mit $f'(x)=3x^2$ lautet der Newton-Schritt
\begin{equation*}
  x_{k+1}=x_k-\frac{f(x_k)}{f'(x_k)}
  =x_k-\frac{x_k^3-3}{3x_k^2}.
\end{equation*}
Gestartet bei $x_0=3$ werden $11$ Iterationen durchlaufen; jede Zeile
\texttt{print(x)} gibt den aktuellen Wert aus. Die Folge konvergiert
quadratisch gegen den Fixpunkt $\sqrt[3]{3}\approx1{,}4422$.

\begin{ergebnis}
a) $\displaystyle\sum_{k=1}^{n}\frac{1}{k^p}$. \\
b) Newton-Iteration
$x_{k+1}=x_k-\frac{x_k^3-3}{3x_k^2}$ für $f(x)=x^3-3$, Start $x_0=3$,
$11$ Schritte, Konvergenz gegen $\sqrt[3]{3}$.
\end{ergebnis}

\begin{kontrolle}
\textbf{Startwerte:} $x_0=3$; $x_1=3-\frac{27-3}{27}=3-0{,}8889=2{,}1111$;
$x_2\approx1{,}6667$; rasch gegen $\sqrt[3]{3}=1{,}44225$. Die Werte folgen
aus der Formel und stimmen mit der Ausgabe überein.

\textbf{Probe:} $\sqrt[3]{3}^{\,3}\approx3{,}0000$ -- bestätigt.
\end{kontrolle}

\begin{merksatz}
\texttt{sum([1/(k**p) for k in range(1,n+1)])} übersetzt zu
$\sum_{k=1}^{n}1/k^p$. Die Schleife wendet zu $f(x)=x^3-3$ die
Newton-Formel an; \texttt{print(x)} zeigt jeden Zwischenstand, die Folge
kriecht gegen die gesuchte Kubikwurzel.
\end{merksatz}
